% ============================================================================
% Section I: Introduction
% ============================================================================

\section{Introduction}
\label{sec:introduction}

The global energy transition toward low-carbon power systems demands
sophisticated planning tools that can co-optimize generation, transmission,
and storage investments while capturing the operational complexity of
variable renewable energy sources and hydro cascades.  Generation and
Transmission Expansion Planning (GTEP) is the class of optimization
problems that determines the minimum-cost combination of infrastructure
investments (CAPEX) and dispatch decisions (OPEX) over long planning horizons
under uncertainty~\cite{hemmati2013comprehensive, mahdavi2019tep_survey}.

Traditional GTEP tools---both commercial (PLEXOS, PLP/PROMAX) and
academic---have been implemented in Fortran or commercial modeling
languages, limiting accessibility and extensibility.  The recent emergence
of open-source alternatives such as PyPSA~\cite{brown2018pypsa},
GenX~\cite{jenkins2017genx}, and IDAES-GTEP~\cite{idaes_gtep2024} has
democratized access to planning tools, primarily through Python and Julia
implementations.  However, these interpreted-language tools can face
computational bottlenecks when assembling large sparse LP/MIP problems with
millions of variables and constraints~\cite{frank2016or_opf}.

This paper presents \gtopt{}, a high-performance, open-source GTEP tool
implemented in modern C++26 that bridges the gap between the computational
efficiency of compiled-language solvers and the accessibility of
open-source, JSON/Parquet-based data interfaces.  \gtopt{} descends from
the FESOP (Fabulous Energy System Optimizer) framework~\cite{buitrago2022fesop},
which extended long-term hydrothermal scheduling tools~\cite{pereira2016hydrothermal}
to include capacity expansion and renewable integration.

The main contributions of this work are:

\begin{enumerate}
  \item A complete mathematical formulation of the GTEP problem as a sparse
    LP/MIP, including DC optimal power flow, battery storage, hydro cascades,
    spinning reserves, and multi-stage capacity expansion
    (Section~\ref{sec:formulation}).

  \item Three solution methods---monolithic LP, SDDP with multi-scenario
    apertures, and a novel multi-level cascade decomposition---with formal
    equivalence guarantees (Section~\ref{sec:methods}).

  \item A modular software architecture with pluggable solver backends,
    native Apache Parquet I/O, and companion tools for data preparation
    (Section~\ref{sec:architecture}).

  \item Validation against IEEE benchmark networks (4-bus through 57-bus)
    including battery storage and capacity expansion variants, with
    cross-validation against pandapower DC OPF results
    (Section~\ref{sec:validation}).
\end{enumerate}

The remainder of this paper is organized as follows.
Section~\ref{sec:literature} reviews the state of the art in GTEP tools
and methods.  Section~\ref{sec:formulation} presents the mathematical
formulation.  Section~\ref{sec:methods} describes the three solution
methods.  Section~\ref{sec:architecture} details the software architecture.
Section~\ref{sec:validation} presents validation results on IEEE benchmark
networks.  Section~\ref{sec:conclusions} concludes with a summary and
future research directions.
